\documentclass[11pt,a4paper]{article}
\usepackage[T1]{fontenc}
\usepackage[utf8]{inputenc}
\usepackage{lmodern}
\usepackage{amsmath,amssymb}
\usepackage[margin=23mm]{geometry}
\usepackage{microtype}
\setlength{\parindent}{0pt}
\setlength{\parskip}{0.5em}
\title{Abatement coefficient: minimal correction\\to the handwritten calculation}
\author{}
\date{}
\begin{document}
\maketitle
\vspace{-2em}

\textbf{The only change is the elastic reference strain.}
Retain the same algebra and notation, but anchor the elastic reference
line at the point on the monotonic curve corresponding to $\sigma_P$:
\[
\varepsilon_P=\frac{\sigma_P}{E}+\varepsilon_{M,p}(\sigma_P),
\qquad
\boxed{\varepsilon_{\mathrm{el}}
=\varepsilon_P+\frac{\sigma_{\mathrm{el}}-\sigma_P}{E}
=\frac{\sigma_{\mathrm{el}}}{E}+\varepsilon_{M,p}(\sigma_P).}
\]
Thus $\varepsilon_{\mathrm{el}}$ is no longer the strain on the
origin-based Hooke line. Here $\sigma_{\mathrm{el}}$ denotes the elastic
reference stress associated with this shifted construction; it must not
be identified automatically with the original elastic-analysis stress.

\textbf{1. Start with the same follow-up relation:}
\begin{align*}
r_M&=-\frac{E(\varepsilon-\varepsilon_{\mathrm{el}})}
                  {\sigma-\sigma_{\mathrm{el}}},\\
\varepsilon&=\frac{\sigma}{E}+\varepsilon_{M,p}(\sigma),
\qquad
\varepsilon_{\mathrm{el}}=\frac{\sigma_{\mathrm{el}}}{E}
                         +\varepsilon_{M,p}(\sigma_P).
\end{align*}
Substitution gives
\begin{align*}
r_M
&=-\frac{\sigma-\sigma_{\mathrm{el}}
 +E[\varepsilon_{M,p}(\sigma)-\varepsilon_{M,p}(\sigma_P)]}
 {\sigma-\sigma_{\mathrm{el}}},\\
r_M+1
&=-\frac{E[\varepsilon_{M,p}(\sigma)-\varepsilon_{M,p}(\sigma_P)]}
 {\sigma-\sigma_{\mathrm{el}}}
=\frac{E[\varepsilon_{M,p}(\sigma)-\varepsilon_{M,p}(\sigma_P)]}
 {\sigma_{\mathrm{el}}-\sigma}.
\end{align*}

\textbf{2. Keep the same definition of the reduction ratio:}
\[
g=\frac{\sigma-\sigma_P}{\sigma_{\mathrm{el}}-\sigma_P}
\qquad(\sigma>\sigma_P).
\]
Multiply numerator and denominator by $1+r_M$ and split the denominator,
exactly as in the handwritten calculation:
\begin{align*}
g
&=\frac{(\sigma-\sigma_P)(1+r_M)}
        {(\sigma_{\mathrm{el}}-\sigma_P)(1+r_M)}\\[0.4em]
&=\frac{(\sigma-\sigma_P)(1+r_M)}
 {(\sigma_{\mathrm{el}}-\sigma)(1+r_M)
  +(\sigma-\sigma_P)(1+r_M)}\\[0.4em]
&=\boxed{\frac{(\sigma-\sigma_P)(1+r_M)}
 {E[\varepsilon_{M,p}(\sigma)-\varepsilon_{M,p}(\sigma_P)]
  +(\sigma-\sigma_P)(1+r_M)}}.
\end{align*}

\textbf{Scope.}
This reproduces the RCC-MRx expression by a pressure-referenced scalar
construction. Using the prescribed code factor $r_M$ in this construction
is an assumption, not a result of the algebra. It is not a general
derivation from the original, origin-based follow-up definition.
If $\sigma_P=0$ and $\varepsilon_{M,p}(0)=0$, the reference shift vanishes
and the calculation reduces to the original one.

\end{document}
